% Section V -- Discussion (condensed). Every number appears in Section IV.

\section{Discussion}

\subsection{Two Costs, Two Remedies}

The energy cost of differential privacy separates into a fixed per-round
overhead and a convergence penalty, and the two call for different responses.
The overhead is an implementation cost --- per-sample gradients and clipping
--- and is what a systems contribution can attack, through better kernels or
amortised clipping. The convergence penalty is bounded by the privacy
accounting rather than by any implementation, and can only be answered through
the round budget.

An evaluation reporting energy per round sees only the first and concludes
that privacy costs a flat half again; measured against the accuracy a
deployment needs, the second is several times larger. This is why the third
axis is worth adding: it does not merely rank configurations, it changes which
engineering effort is worth spending.

\subsection{Privacy Sets a Quality Ceiling}

The cost of privacy is not a constant multiplier. It is invisible at $0.40$
accuracy on HAR, $5.8\times$ at $0.70$, and unbounded at $0.85$. A budget
therefore implies a maximum achievable quality, past which the trade is not
energy for accuracy but nothing for nothing.

This reframes budget selection. The usual formulation --- choose
$\varepsilon$, then see what utility follows --- treats the two as sequential.
Our results suggest choosing them jointly: a deployment needing $0.85$
activity recognition accuracy cannot have $\varepsilon = 1$ under this
accounting, and discovering that after provisioning hardware is expensive. The
energy axis is what makes the ceiling visible, because a utility-only
evaluation reports a lower number without indicating whether more computation
would fix it.

The reason is that the round budget is part of the privacy budget. Every round
a client participates in consumes budget, so $\sigma$ is fixed before training
by a decision federated learning usually treats as free. On HAR, extending
from 30 to 100 rounds --- which the non-private model needs --- raises $\sigma$
at $\varepsilon = 0.5$ from $10.2$ to $18.4$. Training to convergence is
self-defeating: the extension is what prevents convergence. The stopping rule
follows, and at $\varepsilon = 0.5$ on HAR $83\%$ of a run's energy falls
after the peak. We state it conditionally, since HAR reproduces it only at its
tightest budget while CIFAR-10 reproduces it across four; what generalises is
the direction, not the ordering.

\subsection{When the Model Is Noise-Limited, Tuning Is Not the Lever}

Our HAR ablations are negative and useful. At $\varepsilon = 1$ neither local
epochs nor the clipping norm moves accuracy by more than $3\%$ against a
$9.46\%$ floor. A practitioner responding to a disappointing private model by sweeping
hyperparameters will spend considerable energy to land within noise of where
they started; the levers that remain are the privacy and round budgets.

This does not hold in every regime, and CIFAR-10 shows why: there the private
model retains more signal, the same settings matter, and they can be got badly
wrong. Increasing local work per round is the standard way to reduce
communication in federated learning; under DP it stops paying on HAR and
actively harms on CIFAR-10, because local steps enter the budget and force a
larger $\sigma$. Where the boundary between those regimes falls is worth
characterising, and we have only bracketed it.

\subsection{Where the Computation Runs}

Because the cost is measured in joules, it inherits the carbon intensity of
wherever they are drawn: an $18.8\times$ range between Algeria's grid and
Norway's for identical computation and an identical guarantee. Two deployments
with the same regulatory requirements and the same model carry environmental
costs an order of magnitude apart.

The point is sharper for federated learning than for centralised training: the
computation is distributed across wherever the participating devices are, and
no operator chooses the grid. A deployment cannot relocate to a cleaner one
the way a data centre can; what it can adjust is how much computation the
privacy requirement demands.

\subsection{Measuring Energy Around a Privacy Accountant}

Instrumentation attached to a privacy mechanism tends to scale with the
privacy parameter, so its cost mimics the effect under study. The accountant's
noise-calibration search is the case we hit: run inside the federated loop it
reproduced our entire observed $\varepsilon$ trend while executed step counts
stayed flat. Nothing in the resulting numbers announces this.

We recommend three controls for future work on the energy of privacy
mechanisms: derive noise multipliers before measurement begins, verify per
round that executed and calibrated step counts agree, and randomise run order
so that drift over a sweep cannot be mistaken for an $\varepsilon$ effect.

\subsection{Threats to Validity}

Clients are virtual and share one GPU, so the ratios transfer to real devices
but the absolute joules do not; confirming them on physical hardware is the
immediate next step. Measurement is GPU-only and communication energy is
modelled --- for a $0.15$--$0.24$~MB model over tens of rounds we expect it
small, but have not shown it. Two datasets bound the effects rather than
establishing them. And our accounting claims no amplification by subsampling,
so the reported $\varepsilon$ is an upper bound: these are the costs under the
weakest trust assumptions, not the best achievable.
